%20 min preso!
\documentclass[xcolor=table]{beamer}
\usepackage{beamerthemesplit}
\usepackage{wrapfig}
\usetheme{SPbGU}
\usepackage{pdfpages}
\usepackage{amsmath}
\usepackage{cmap}
\usepackage[T2A]{fontenc}
\usepackage[utf8]{inputenc}
\usepackage[english]{babel}
\usepackage{indentfirst}
\usepackage{amsmath}
\usepackage{tikz}
\usepackage{multirow}
\usepackage[noend]{algpseudocode}
\usepackage{algorithm}
\usepackage{algorithmicx}
\usepackage{fancyvrb}
\usetikzlibrary{calc}
\usetikzlibrary{shapes,arrows}
\usetikzlibrary{arrows,automata}
\usetikzlibrary{positioning}

\usepackage{tabularx}
\newcolumntype{Y}{>{\raggedleft\arraybackslash}X}


\newtheorem{mytheorem}{Theorem}
\renewcommand{\thealgorithm}{}

\newcommand{\tikzmark}[1]{\tikz[overlay,remember picture] \node (#1) {};}
\def\Put(#1,#2)#3{\leavevmode\makebox(0,0){\put(#1,#2){#3}}}

\tikzset{
    state/.style={
           rectangle,
           rounded corners,
           draw=black, very thick,
           minimum height=2em,
           inner sep=2pt,
           text centered,
           },
}

\beamertemplatenavigationsymbolsempty

\title[Parsing Techniques for CFPQ]{Parsing Techniques for Contex-Free Path Querying}
%\subtitle[YaccConstructor]{Parsing techniques for graph analysis}
% То, что в квадратных скобках, отображается в левом нижнем углу.
\institute[JetBrains Research]{
JetBrains Research, Programming Languages and Tools Lab  \\
Saint Petersburg University
}

% То, что в квадратных скобках, отображается в левом нижнем углу.
\author[Semyon Grigorev]{\textbf{Semyon Grigorev}}

\date{April 05, 2019}

\begin{document}
{
\begin{frame}[fragile]
  \begin{table}
  \centering
  \begin{tabularx}{\linewidth}{YcX}
    \includegraphics[height=1.5cm]{pictures/jetbrainsResearch.pdf} \hfill
    & \begin{minipage}[t]{0.3\textwidth}\center \vspace{-1cm}
      \end{minipage}
    & \hfill \includegraphics[height=1.5cm]{pictures/SPbGU_Logo.png}
  \end{tabularx}
  \end{table}
  \titlepage
\end{frame}
}

\begin{frame} \frametitle{Programming Languages and Tools Lab}
    \begin{itemize}
      \item \url{https://research.jetbrains.org/groups/plt_lab}
    \end{itemize}
\end{frame}


\begin{frame} \frametitle{Formal languages for data analysis}

\begin{itemize}
  \item Semyon Grigorev
\end{itemize}

\end{frame}

\begin{frame} \frametitle{Topics of interest}
\begin{itemize}
\item Formal language theory
\item Parsing algorithms
\end{itemize}
\end{frame}


\begin{frame} \frametitle{Formal language constrained path querying}
\begin{itemize}
\item Finite directed edge-laballed graph $\mathcal{G} = (V,E,L)$
\item The path is a world over $L$: $\omega(p) = \omega(v_0 \xrightarrow{l_0} v_1 \xrightarrow{l_1} \dots \xrightarrow{l_{n-1}} v_n ) = l_0 \cdot l_1 \cdot \ldots \cdot l_{n-1}$
\item The language $\mathcal{L}$ (over $L$)
\end{itemize}
\pause
\begin{itemize}
  \item Reachability problem: $Q=\{(v_i,v_j) \ | \ \exists p = v_i \dots v_j, \omega(p) \in \mathcal{L}\}$
  \item Path querying problem: $Q=\{p \ | \ \omega(p) \in \mathcal{L}\}$
  \begin{itemize}
    \item Single path, all paths, shortest path \dots
  \end{itemize}
\end{itemize}

\end{frame}

\begin{frame} \frametitle{Context-Free path querying}
\begin{itemize}
\item $\mathcal{L}$ is a context-free language
\item $G_{\mathcal{L}} = (N,\Sigma,R,S)$
\item Reachability problem: $Q=\{(v_i,v_j) \ | \ \exists p = v_i \dots v_j, S \xrightarrow[G_L]{*} \omega(p) \}$
\item Path querying problem: $Q=\{p \ | \ \omega(p) \in \mathcal{L}\}$
\end{itemize}

\end{frame}


\begin{frame} \frametitle{Example of CFPQ}
\begin{itemize}
\item !!!!
\end{itemize}
\end{frame}


\begin{frame} \frametitle{Applicability}
\begin{itemize}
\item Static code analysis
\item Graph data bases
\item !!!!
\end{itemize}
\end{frame}


\begin{frame} \frametitle{Open questions}
\begin{itemize}
\item Subcubic (or to BMM)
\item Engineering
\item New field for parsing algorithms
\end{itemize}
\end{frame}


\begin{frame} \frametitle{!!!!!}
\begin{itemize}
\item Subcubic (or to BMM)
\item Engineering
\item New field for parsing algorithms
\end{itemize}
\end{frame}


\begin{frame} \frametitle{Bioinformatics. Sequences analysis}
\begin{itemize}
\item High performance
\item New classes of grammars
\item !!!!
\end{itemize}
\end{frame}


\begin{frame}
\frametitle{Contact Information}
\begin{itemize}
  \item Semyon Grigorev:
    \begin{itemize}
      \item \href{mailto:s.v.grigoriev@spbu.ru}{s.v.grigoriev@spbu.ru}
      \item \href{mailto:Semen.Grigorev@jetbrains.com}{Semen.Grigorev@jetbrains.com}
    \end{itemize}
  \item Polina Lunina:
  \begin{itemize}
    \item  \href{mailto:lunina_polina@mail.ru}{lunina\_polina@mail.ru}
  \end{itemize}
  \item Trained models: \href{https://github.com/YaccConstructor/YC.Bio}{https://github.com/YaccConstructor/YC.Bio}
\end{itemize}
\vspace{2cm}
\center{\huge{Thanks!}}
\end{frame}

\end{document}
